\chapter{Introdução}
\label{cap:introducao}

As redes elétricas de distribuição são tradicionalmente passivas, ou seja, 
a potência flui de forma unidirecional da geração para o consumo~\cite{grainger1994}. 
Nesse modelo, a qualidade da energia, em especial a regulação de tensão e o fator de potência, 
é assegurada pela robustez da geração centralizada, como hidrelétricas e termelétricas, 
pouco sensível a variações de curto prazo. A crescente integração de fontes de geração distribuída 
(GD), como os sistemas fotovoltaicos (FV), transformou essas redes em redes ativas, nas quais o 
fluxo de potência passa a ser bidirecional, determinado tanto pela geração quanto pela carga. 
Tal transformação impõe novos desafios à operação e ao controle, sobretudo no que tange à regulação 
de tensão.

A disseminação da GD criou demanda por Sistemas de Armazenamento de Energia em Baterias 
(\emph{Battery Energy Storage Systems} --- BESS), empregados para armazenar o excedente de energia 
em horários de baixa demanda e injetá-lo na rede nos horários de ponta. A eficácia desses sistemas, 
contudo, depende diretamente da estratégia de controle adotada. O controle centralizado, 
tradicionalmente utilizado, apresenta limitações de escalabilidade e robustez, pois exige infraestrutura 
de comunicação complexa e é suscetível a falhas em um único ponto de processamento. Em contrapartida, 
as estratégias de controle distribuído permitem que cada unidade BESS tome decisões com base em 
informações locais e na comunicação parcial com vizinhos imediatos~\cite{nabatirad2022}, aumentando a 
confiabilidade e facilitando a operação \emph{plug-and-play} de novos dispositivos.

Os conjuntos formados por GD e BESS favorecem um controle distribuído e automático da rede por meio dos 
inversores que conectam tanto os geradores fotovoltaicos quanto os BESS. Ao controlar a potência ativa e 
reativa injetada ou absorvida, é possível regular a tensão localmente. À medida que cresce o número de 
BESS autônomos e distribuídos, torna-se necessário coordená-los, de modo que operem de forma harmônica 
também em nível global, sem deteriorar os perfis de tensão e o carregamento agregado da rede.

\section{Motivação e definição do problema}
\label{sec:motivacao}

Na literatura, há três grandes classes de soluções para o controle de BESS em redes de distribuição. A 
primeira é o despacho ótimo~\cite{chreim2024,li2022}, baseado na otimização de um problema de programação 
linear inteira mista (\emph{Mixed-Integer Linear Programming}) que calcula o perfil de carga e descarga 
do BESS para cada intervalo de tempo a partir de previsões de carga e geração; trata-se de uma estratégia 
em malha aberta, na qual desvios entre previsão e realidade degradam o desempenho. A segunda é o controle 
centralizado, em que dados de toda a rede são coletados para gerar ações em um único ponto de controle, 
exigindo infraestrutura de comunicação robusta e apresentando ponto único de falha. A terceira é o 
controle distribuído em dupla camada, no qual cada BESS age localmente de forma autônoma na primeira 
camada e se coordena com os vizinhos na segunda.

Apesar do potencial das soluções distribuídas, persistem lacunas relevantes. A maior parte dos trabalhos 
não analisa como o desequilíbrio de fases, característico das redes reais, afeta o algoritmo de controle, 
e poucos incorporam modelos de degradação e ciclo de vida das baterias ao laço de controle, aspecto 
diretamente ligado à viabilidade econômica do uso de BESS para suporte de rede~\cite{gangwar2024}. Este 
projeto busca validar e estender a solução proposta por~\cite{paradacontzen2024}, implementando, nesta 
primeira etapa, uma estratégia análoga de controle distribuído em dupla camada e avaliando seu desempenho 
de regulação de tensão em uma rede de teste padrão.

\section{Objetivos}
\label{sec:objetivos}

O objetivo geral deste trabalho é desenvolver e avaliar uma estratégia de controle distribuído em dupla 
camada para BESS autônomos, voltada à regulação de tensão em redes de distribuição de média tensão com 
alta penetração de geração fotovoltaica. No escopo desta primeira etapa do Projeto de Formatura, 
definem-se os seguintes objetivos específicos:

\begin{itemize}
  \item revisar a literatura de estratégias de controle distribuído e de algoritmos de consenso aplicados à coordenação de BESS;
  \item desenvolver um ambiente de cossimulação acoplando Python e o software OpenDSS para a execução de fluxos de potência iterativos em séries temporais;
  \item elaborar cenários de penetração de recursos energéticos distribuídos (FV e BESS) por meio do método de Monte Carlo;
  \item implementar o controle primário (local) e o controle secundário (colaborativo, por consenso) sobre as unidades distribuídas;
  \item avaliar o desempenho da estratégia em termos de regulação de tensão, comparando os casos sem controle, com controle primário isolado e com as duas camadas, em uma rede de teste padrão da IEEE.
\end{itemize}

A extensão da solução a redes desequilibradas, a incorporação dos modelos de degradação e ciclo de vida 
das baterias e a aplicação a uma rede de distribuição real constituem objetivos da segunda etapa do 
projeto e são tratados, neste relatório, apenas como trabalhos futuros (Capítulo~\ref{cap:conclusao}).

\section{Organização do trabalho}
\label{sec:organizacao}

O restante deste relatório está organizado da seguinte forma. O Capítulo~\ref{cap:fundamentacao} 
apresenta a fundamentação teórica e a revisão bibliográfica das estratégias de controle distribuído e 
dos algoritmos de consenso. O Capítulo~\ref{cap:ambiente} descreve o ambiente de cossimulação 
Python--OpenDSS e as redes de teste utilizadas. O Capítulo~\ref{cap:cenarios} detalha a geração de 
cenários pelo método de Monte Carlo. O Capítulo~\ref{cap:controle} formaliza a estratégia de controle em 
duas camadas. O Capítulo~\ref{cap:resultados} apresenta e discute os resultados. Por fim, o 
Capítulo~\ref{cap:conclusao} sintetiza as conclusões, mapeando-as aos objetivos, e aponta os próximos 
passos.
